% ============================================================
% s03-chI-3.tex —— Chapter I 块 C：§I.10（承接，p.39）--§I.13（p.44）
% 底本：arXiv:2102.13459v4，PDF p.39--p.44
% 说明：§I.10「谱作用」的 \section 与小节前两条编号对象（Thm I.10.1、
%       Conj I.10.2）在块 B 的 s02-chI-2.tex 中，本文件不自开 I.10 节，
%       以免整书节号与定理号错位；本文件从 p.39 上的 Remark I.10.3 续起。
% ============================================================

\begin{remark}
考虑系数取在 $\ol{\QQ}_\ell$ 中的猜想。\srcpage{39}理想地，该猜想还应包含 $t$-结构之间的比较。遗憾的是，我们当时没有立即看到一个可用于匹配 $t$-结构的合适候选。理想地，这会把左侧的反常 $t$-结构（它由抽象的理由即已良定义，并至少隐含地出现在 [CS17]、[CS19a] 中；对局部—整体相容性问题而言，它似乎正是“正确”的 $t$-结构）与右侧某个“反常—凝聚”（perverse-coherent）$t$-结构相比较。\yizhu{“反常—凝聚”（perverse-coherent）$t$-结构：glossary 仅收 perverse t-structure＝反常 $t$-结构；原文给该词加引号，说明尚未定名，译文保留英文原词。}{A：原文 §I.10 p.39；glossary.tsv}若果真如此，该等价还将给出两边阿贝尔心之间不可约对象的一个双射。在左端，这些不可约对象将可由 $B(G)$ 中元素 $b$ 与 $G_b(E)$ 的不可约光滑表示 $\pi_b$ 组成的对 $(b,\pi_b)$，借助中间扩张来枚举。在右端，它们很可能对应于一个 Frobenius-半单的 $L$-参数 $\varphi:W_E\to\hat G(\ol{\QQ}_\ell)$ 连同 $\varphi$ 的中心化子 $S_\varphi$ 的一个不可约表示。独立于这个范畴猜想，人们可以问：这两个集合事实上是否典范地双射？\footnote{此问题已由 Bertoloni Meli--Oi [BMO22] 肯定地回答。Hansen [Han24] 更进一步，通过引入 $\Bun_G$ 上的 hadal $t$-结构，推进了对可能匹配的 $t$-结构的理解。}\yizhu{hadal 指海洋的“深渊带”，此处是作者对 Hansen 所引入 $t$-结构的戏称；文献标签 [Han24] 照录不译。}{A：原文 §I.10 p.39 脚注 7；Hansen [Han24]}
\end{remark}

\section{思想的来源}

最后，让我们从自己有偏的视角，对这些思想的来龙去脉作一番交代。让我们先回顾自己早期在局部 Langlands 对应方向上的一些工作。Fargues [Far04] 证明了：对于 $\mathrm{GL}_n$（以及 $U(3)$）的基本 Rapoport--Zink 空间与一般的 minuscule 余特征标，局部 Langlands 对应的一个适当版本在上同调中得到了实现。此外，Fargues [Far08] 还证明了 Lubin--Tate 塔与 Drinfeld 塔之间的对偶同构。此时 Fargues 已把它看作将 Jacquet--Langlands 对应几何化的一次尝试，见 [Far08, Theorem 2 of the Préambule]。另一方面，Scholze [Sch13] 给出了 $\mathrm{GL}_n$ 的局部 Langlands 对应的一个新证明。他的结果指向这样一个想法：在 $p$-可除群的模叠上应当存在某些层（当限制到完满概形上时，它们可视为上文考虑过的叠 $\mathrm{GL}_n\text{-}\Isoc$ 的“一部分”），从而给出局部 Langlands 对应的某种几何化，后者随即被表述为某种特征标层性质（受 [Sch13] 中特征标公式的启发）。Boyer（例如见 [Boy09]）以及 Dat 的未发表工作中也有相关的观察。然而，Scholze 一直对 $p$-可除群叠极差的几何性质感到不安。

此时，我们两人基本上都已把局部 Langlands 抛在脑后，转而研究别的问题。Fargues 在与 Fontaine 的合作 [FF18] 中发现了 $p$-进 Hodge 理论的基本曲线；对 Fargues 而言，一个最初的关键动机是发展“不带 Galois 作用的 $p$-进 Hodge 理论”，即针对像 $\mathbb{C}_p$ 这样的域的理论。事实上，他在关于 Rapoport--Zink 空间的某些工作中正需要这个。另一方面，Scholze 发展了完满空间 [Sch12]，其动机来自权—单值猜想。2011 年 3 月，在普林斯顿的一次会议上，Weinstein 在其报告之后又作了一场关于他在无穷层 Lubin--Tate 塔上结果的报告，这使人们清楚：它事实上是一个完满空间。Scholze 当时已很想理解 Lubin--Tate 塔与 Drinfeld 塔之间的同构，而现在清楚了，它本当是完满空间之间的同构。这在 [SW13] 中完成。写作 [SW13] 时，Fargues--Fontaine 曲线的视角已变得中心，我们意识到，两塔之间的同构不过就是 Fargues--Fontaine 曲线上极小\srcpageinl{40}修正 $\cO_X^n\to\cO_X(\frac{1}{n})$ 空间的两种对偶描述，取决于固定的是哪个丛、作修正的又是哪个丛。这是局部 Langlands（以 Lubin--Tate 空间与 Drinfeld 空间的上同调为其编码）与 Fargues--Fontaine 曲线上向量丛理论之间的第一个清晰联系，尽管 Scholze 当时并未足够认真地对待它。此外，Fargues 在 [FF18] 中证明“弱可容许蕴含可容许”时已经注意到，向量丛的修正扮演着重要角色：一个滤过 $\varphi$-模的 Hodge 滤过使人能够修改与某个等晶体（isocrystal）相伴的向量丛，从而定义一个新的向量丛——也就是他本人在纪念 Jean-Marc Fontaine 的会议报告 [Far10] 中所说的“施行一次 Hecke 对应”。

这种对偶视角也把两个对偶的周期态射推到了关注的中心：Hodge--de Rham 周期映射，与 Hodge--Tate 周期映射（二者在对偶同构下互换）。把 Lubin--Tate 塔看作椭圆曲线模空间的一部分，Scholze 随即意识到 Hodge--Tate 周期映射甚至在无穷层椭圆曲线模空间上整体存在（在 Berkovich 拓扑空间的层面上，Fargues 此前也已观察到这一点）。此外，Scholze 意识到 Hodge--Tate 周期映射给出了从椭圆曲线模空间到 $p$-可除群模空间之映射的一个替代物，并且他为 [Sch13] 的几何解释所寻求的那些层，更有可能存在于 Hodge--Tate 周期映射的靶上——后者不过是 $\mathbb{C}_p$ 上的一个射影空间；他在一次 MSRI 报告 [Sch14] 中勾勒了这些想法。（同样地，Dat 也有类似的想法。）最终，这一视角被用在他与 Caraiani 的合作 [CS17]、[CS19a] 中，以研究 Shimura 簇上同调中的挠。与 Caraiani 的合作需要 Fargues--Fontaine 曲线上 $G$-主丛的分类，而这已由 Fargues [Far20a] 证明。

越来越多地采取“只让完满空间映入”来研究一切几何对象的视角，菱形的想法便迅速浮现，其中包括得到若干份 $\Spec\QQ_p$ 的可能性（这一想法最早见于已发表文献的形态是 [Wei17]），以及定义一般的 $p$-进 shtuka 模空间。这些想法在 2014 年秋 MSRI 学期期间 Scholze 的 Berkeley 课程 [SW20] 中得到阐述。最终的目标始终是把 V. Lafforgue 的工作 [Laf18] 适配到 $p$-进域的情形；最初的策略是经由局部 shtuka 模空间的上同调来定义所需的游走算子。学期之初，Scholze 对几何 Langlands 纲领仍非常戒慎，因为它似乎无法纳入 $p$-进群超尖（supercuspidal）表示那些微妙的算术性质。于是，下面这一点完全是一个出人意料的观念飞跃：对整个主题而言最好的视角，其实是在 Fargues--Fontaine 曲线上把局部 Langlands 对应看作一个几何 Langlands 对应——这是 Fargues 在 MSRI 的一次茶歇上提出的（部分受启发于他就曲线上 $G$-丛空间与 [Far20a] 的关系所作的密集思考）。Fargues 当时采取的是 Hecke 特征层的视角，想要对任意（离散）$L$-参数 $\varphi$ 构造出 $\Bun_G$ 上一个以 $\varphi$ 为特征值的相伴 Hecke 特征层 $A_\varphi$。这应当定义一个函子 $\varphi\mapsto A_\varphi$，从而带上 $\varphi$ 的中心化子群 $S_\varphi\subset\hat G$ 的一个作用，而 $A_\varphi$ 相应的 $S_\varphi$-等型分解应当实现 $L$-包的内部结构。此外，Hecke 特征层性质应当蕴含关于局部 Shimura 簇上同调的 Kottwitz 猜想 [RV14, Conjecture 7.3]。这让一切都融贯起来了。特别地，它为 $L$-包的内部结构给出了一个令人信服的几何起源，并且也与 Kaletha 晚近的工作 [Kal14] 相合，后者用基本的 $G$-等晶体来精细研究 $L$-包。

\srcpage{41}
遗憾的是，这个猜想是在极其不牢靠的基础上表述出来的：它假定人们可以把模叠 $\Bun_G$ 当作通常代数几何的对象来操作。当然，它还假定存在几何 Satake 的一个版本，等等（原文作“etc.pp.”）。\yizhu{原文作“etc.pp.”，为拉丁语套语 “et cetera perge perge”（等等、不一而足）的缩写，非笔误；译文按“等等”处理。}{A：原文 §I.11 p.41}另一方面，我们意识到，一旦我们能够仅仅把 Fargues 的猜想表述出来，就有足够的机器可用，以把 Lafforgue 的想法 [Laf18] 应用于得到“自守→Galois”方向，并定义（半单）$L$-参数（正如 Genestier--Lafforgue [GL17] 在等特征情形所作的那样）。

自那以后，这是一段漫长且非常痛苦的过程。第一步是为与这幅图景相关的几何对象范畴——即菱形——给出一个好的定义。特别地，人们必须证明相关的仿射 Grassmann 空间具有这一性质。这是 Berkeley 课程 [SW20] 的主要结果。为作证明，引入了 v-层的概念，它此后在代数几何中也自成一体地发展起来（参看 [BM21]）。（一般地，v-下降被证明是一种极有力的证明技术。我们在此用它重新证明关于 Fargues--Fontaine 曲线的基本定理，毫不费力地回收了 [FF18] 与 [KL15] 的主要定理。）接下来，必须为菱形的 \'etale 上同调发展一套六函子形式化，这已在 [Sch17a] 中实现，至少对挠系数而言。过渡到 $\QQ_\ell$-系数比概形情形费力更多，我们将在下文评述。[Sch17a] 的一项中心技术是 pro-\'etale 下降，更一般地是 v-下降。事实上，[Sch17a] 的几乎所有定理都是用这类下降技术证明的，本质上把它们归约为几何点的射有限族。令 Scholze 惊讶的是，这种把光滑空间拆解为射有限集合的过程竟然在证明几何结果上有力量，而这一认识极大地推动了凝聚数学的发展（后者反过来又回馈到本项目之中）。

至此，才有可能去思量 Fargues 的猜想。在这方面，首先要建立的结果是 $D_{\text{ét}}(\Bun_G,\ZZ/\ell^n\ZZ)$ 性态良好，例如对“可容许”层满足 Verdier 双对偶。我们找到了一个证明，它以 $\Bun_G$ 在任意 $b\in B(G)$ 附近的某个“图卡” $\pi_b:\cM_b\to\Bun_G$ 的上同调光滑性为条件；这在 Scholze 的 IH\'ES 课程 [Sch17b] 中有解释。虽然对于 $G=\mathrm{GL}_n$，$\pi_b$ 的上同调光滑性可以通过直接强攻来证明，但在一般情况下，我们只能把它表述为下述一般“光滑性的 Jacobian 判别法”的一个特例：该判别法针对的是参数化某个 Fargues--Fontaine 曲线上光滑 adic 空间 $Z$ 的截面 $Z\to X_S$ 的空间。证明这一 Jacobian 判别法需要另外三个关键想法。第一个是“形式光滑”的概念，其中到无穷小加厚（在完满几何中并不存在）的提升被替换为到实际的小开（或 \'etale）邻域的提升。由此得到的概念与经典拓扑中绝对邻域收缩核的概念密切相关 [Bor67]。通过一番真正的“分析”，不难证明截面空间是形式光滑的。遗憾的是，这似乎还不足以保证上同调光滑性。第一个问题是，形式光滑并不蕴含任何有限维性。这里就出现了第二个关键想法，即 Bhatt 的洞见 [BS22]：Zariski 闭浸入在 [Sch15, Section II.2] 的意义下是强 Zariski 闭的（与 Scholze 在那里所作的断言相反）。至此，只需证明：在一个有限维完满球中既形式光滑又 Zariski 闭的空间是上同调光滑的。遗憾的是，尽管屡经尝试，我们仍然无法证明 [Sch17a] 中的不同维数概念（Krull 维数、dim. trg、上同调维数）对这样的空间彼此一致。这很可能正是该理论中最重要的基础性开问题：

\srcpage{42}
\begin{question}
设 $X\subset\wt{\mathbb{B}}^n_C$ 是 Zariski 闭的，其中 $\wt{\mathbb{B}}^n$ 是一个完满球\yizhu{“完满球”即完满空间意义下的 $n$ 维球；glossary 仅收 perfectoid space＝完满空间，本词按词根一致暂译，待并入术语表。}{A：原文 §I.11 p.42；glossary.tsv}。证明 $X$ 具有性态良好的维数。
\end{question}

事实上，我们觉得这简直不可思议：我们能够证明各式各样的非平凡几何结果，却始终无法毫不含糊地谈论维数！

针对这一缺陷，我们在进攻时用到了第三个关键想法：即普遍局部无环层的概念，我们为了证明几何 Satake 也独立地发展了它。\yizhu{原文此处作 “Our attacks on this failing, a third key idea comes in:”，疑脱一 “In”（当作 “In our attacks on this failing”）；译文按原意补足。}{A：原文 §I.11 p.42}容易看出，形式光滑加上有限维性蕴含常层是普遍局部无环的；剩下的就是看出对偶层是可逆的。这可以通过形变到法锥来证明（利用普遍局部无环性把法锥上的结果传播到邻域）。我们是在 2018 年 7 月 Luminy 的一次会议上找到这个论证的；使用形变到法锥的灵感，可能来自 Clausen 在“线性化猜想”证明中的用法。

这些结果足以证明 $D_{\text{ét}}(\Bun_G,\ZZ/\ell^n\ZZ)$ 性态良好，并且已足以证明关于 Rapoport--Zink 空间上同调（带挠系数）的若干新的有限性结果。我们接下来的重点是几何 Satake。这本质上需要普遍局部无环层理论，以及 Braden 双曲局部化定理 [Bra03] 的一个版本。我们为二者都找到了替代品。关于普遍局部无环层，我们能够证明大多数基本定理的相应版本，然而未能证明它们一般而言在相对 Verdier 对偶下保持（尽管在所有相关情形我们都能手工验证）。Lu--Zheng [LZ22] 随后找到了普遍局部无环层的一个新刻画，使得在相对 Verdier 对偶下的稳定性变得显然。他们的论证立刻可搬运到我们的设定中。最终我们使用了一个略有不同的刻画，但精神上论证仍与他们的相同。关于双曲局部化，我们无法遵循 Braden 依赖漂亮坐标选择的论证。取而代之，我们把所有论证归约到下述（易于证明的）原理：若 $X$ 是一个带 $\GG_m$-作用的（部分紧合）空间，且 $[X/\GG_m]$ 为 qcqs，而 $A\in D_{\text{ét}}([X/\GG_m],\Lambda)$，则 $X$ 的以 $A$ 为系数的部分紧支撑上同调为零。这里的想法是，$\GG_m$-作用把 $X$ 收缩到其一端。此后，几何 Satake 的证明大体沿 [MV07] 的路线进行，尽管论证中有若干改进；特别地，我们给出了到秩 $1$ 群的简单归约，并落实了与对偶群的同构。

利用这些结果，一切原料都已就位，但只有用挠系数才行。人们可以形式地过渡到 $\ell$-进完备层，但这会导致研究 Banach $\QQ_\ell$-向量空间上的表示，这非常不自然。在此期间，Clausen 来到 Bonn，Clausen 与 Scholze 开始发展凝聚数学以及固体模的理论 [CS]。他们意识到，人们也可以在概形或菱形上定义固体 $\ZZ_\ell$-层，而这样做使得研究离散的 $\ol{\QQ}_\ell$- 或 $\QQ_\ell$-向量空间上的表示成为可能，正如所愿。我们在此着手这件事：先在任何小 v-叠上定义固体 $\ZZ_\ell$-层，连同某种五函子形式化（其中用相对同调代替紧支上同调；其右伴随则是拉回，故只有 5 个函子），此后过渡到“lisse-\'etale”层的一个子范畴，以定义所需范畴 $D_{lis}(\Bun_G,\ol{\QQ}_\ell)$，它具有恰好所需的性质。

\emergencystretch=2.5em
在此期间，\srcpageinl{43}几何 Langlands 纲领中有 Nadler--Yun [NY19] 与 Gaitsgory--Kazhdan--\allowbreak Rozenblyum--\allowbreak Varshavsky [GKRV22] 的相关工作，它们蕴含：我们所构造出的范畴结构——$D(\Bun_G,\ol{\QQ}_\ell)$\allowbreak 连同 Hecke 算子的作用——\allowbreak 形式地诱导了 $L$-参数叠上完美复形范畴在 $D(\Bun_G,\ol{\QQ}_\ell)$ 上的一个作用，从而给出基于游走算子的 $L$-参数构造的一个范畴升级。（我们当时知道这件事的某种弱形式，限于椭圆参数时成立；这在 [Sch17b] 的最后一讲中讨论过，依据的是 Anschütz 的一些未发表结果。）在此，我们费力去证明一个带 $\ZZ_\ell$-系数的结果。

\emergencystretch=0pt

\section{致谢}

我们为稿件准备的漫长拖延致歉。早在 2016 年 4 月，一次 Oberwolfach Arbeitsgemeinschaft 便对这些想法作了介绍。所需的关于菱形 \'etale 上同调（带挠系数）的基础，基本上在 2016/17 冬季学期完成，当时 Bonn 的 ARGOS 讨论班研读了演变中的稿件 [Sch17a]，这对我们是极大的帮助。接下来，我们在 IH\'ES 讲授了 Hadamard 课程（2017 年春），[Sch17b]。那时我们基本上已获得足够的结果来完成 $L$-参数的构造；缺失的一环是 Jacobian 判别法。特别地，我们明白了必须用局部图卡 $\cM_b$ 来研究 $\Bun_G$。更详细地说，这些结果在 2017 年夏季学期的 ARGOS 讨论班上讨论过，而几何 Satake 的证明则在 2018 年夏季学期的 ARGOS 讨论班上讨论过。我们衷心感谢这些讨论班的所有参加者，感谢他们研读这些稿件并给出宝贵的反馈。在此期间，2018 年春，我们修订了 Berkeley 讲义 [SW20]，并使用 [Sch17a] 给出了 $p$-进 shtuka 模空间的完整构造。然后我们终于在 2018 年 7 月于 Luminy 的一次会议上找到了 Jacobian 判别法的完整证明。我们早期撰写结果的尝试撞上了这样一个问题：我们想先证明关于维数的基本健全性结果。后来，我们又被与 Dustin Clausen 一起发展凝聚数学带偏了方向，而这最终使本文能够使用 $\QQ_\ell$-系数，终于给出最自然的结果。事实上，我们一直想要把基础发展得使诸如 Hansen--Kaletha--Weinstein 关于 Kottwitz 猜想的论文 [HKW22] 能够建立在一个恰当的基础上，并证明关于一般光滑 $\QQ_\ell$-表示的结果；我们相信我们终于做到了这一点。2020/21 冬季学期，Scholze 在 Bonn 开了一门（线上）课程讲授这些结果，他要感谢所有参加者的宝贵反馈。我们感谢许多与我们讨论过这些结果的数学家，包括 Johannes Anschütz、David Ben-Zvi、Joseph Bernstein、Bhargav Bhatt、Arthur-César le Bras、Ana Caraiani、Dustin Clausen、Pierre Colmez、Jean-François Dat、Vladimir Drinfeld、Jessica Fintzen、Dennis Gaitsgory、Toby Gee、David Hansen、Eugen Hellmann、Alex Ivanov、Tasho Kaletha、Kiran Kedlaya、Dmitry Kubrak、Vincent Lafforgue、Judith Ludwig、Andreas Mihatsch、Sophie Morel、Jan Nekovář、Wiesława Nizioł、Vincent Pilloni、Michael Rapoport、Timo Richarz、Will Sawin、Tony Scholl、Antoine Touzé、Roman Travkin、Wilberd van der Kallen、Yakov Varshavsky、Jared Weinstein、Zhiyou Wu、Xinwen Zhu、Konrad Zou，以及无疑还有许多其他人。特别感谢 David Hansen 对初稿的许多评论。他还观察到，该范畴等价若存在则唯一，这是我们在本稿最初版本中未曾作出的观察。Konrad Zou 观察到，在坏素数处，幂零奇异支撑条件的表述中所用的幂零锥可能行为不正确。此外，我们感谢 Toby Gee、David Helm、Sug Woo Shin，特别是 Zhiyu Zhang\srcpage{44} 对第一版的评论与建议。最后，我们衷心感谢审稿人非常仔细的阅读和许多有用的建议。

在本稿写作期间，Fargues 得到 ANR 资助 ANR-14-CE25 PerCoLaTor 与 ERC Advanced Grant 742608 GeoLocLang 的支持，Scholze 得到 DFG Leibniz Prize 以及 DFG 在 Excellence Strategy 之下的 EXC-2047/1 -- 390685813 的支持。

\section{记号}

在本文的大部分篇幅中，$E$ 表示一个非阿基米德局部域，其剩余域 $\FF_q$ 的特征为 $p>0$；我们固定 $\FF_q$ 的一个代数闭包 $k=\ol{\FF}_q$。于是 $\breve{E}$ 是 $E$ 的完全无分歧扩的完备化，其剩余域为 $k$。我们还固定 $E$ 的一个可分闭包 $\ol{E}$，其绝对 Galois 群为 $\Gamma=\Gal(\ol{E}|E)$，其中包含 Weil 群 $W_E$、惯性子群 $I_E$ 与野惯性子群 $P_E$。字母 $P$ 通常表示 $P_E$ 的开子群，但偶尔也用来表示一个有限 $p$-群（或更一般地，一个阶与 $\ell$ 互素的有限可解群）。

群 $G$ 通常是 $E$ 上的一个约化群；约化群一律假定是连通的。

对任何拓扑空间 $X$，我们用 $\ul{X}$ 记这样的层：它把任何 $S$（在相关的测试范畴中，通常是完满空间）送到从 $|S|$ 到 $X$ 的连续映射。这与从拓扑空间到凝聚集的过渡是同一精神，见 [CS]。我们偶尔使用凝聚的语言，但不使用 [CS] 中的任何非平凡结果。特别地，我们对固体 $\ell$-进层的讨论是自足的。

我们偶尔使用“animated”这一术语，见 [CS]、[\v{C}S19b]。特别地，我们用 anima 一词来指 [Lur09] 中所谓的空间、$\infty$-群胚或同伦型；而对任何环 $A$，animated $A$-代数的 $\infty$-范畴是由单纯 $A$-代数反演弱等价而得到的 $\infty$-范畴。于是，animated $A$-代数在筛余极限下由多项式代数 $A[X_1,\dots,X_n]$ 自由生成。\yizhu{sifted colimit 指以筛范畴（sifted category）为指标范畴的余极限；本词未入 glossary，暂译“筛余极限”并保留英文，待并入术语表。}{A：原文 §I.13 p.44；glossary.tsv}

若 $\mathcal{C}$ 是一个带有群 $G$ 作用的（$\infty$-）范畴，我们记 $\mathcal{C}^{BG}$ 为 $\mathcal{C}$ 中 $G$-等变对象的（$\infty$-）范畴。注意这里的数据实际上是一个函子 $BG\to\mathrm{Cat}_\infty$，而 $\mathcal{C}^{BG}$ 按定义是这个图表的极限。（更惯常的写法是 $\mathcal{C}^G$，但那会导致记号不一致。）此外，我们常把分类叠写作 $\ast/G$ 而不是 $BG$，因为字母 $B$ 也用来表示 Borel 子群（而且我们强烈偏好 $\ast/B$ 而非 $BB$），并且它还出现在 Kottwitz 集 $B(G)$ 中。

\subsection*{记号一览（译者整理）}

\noindent 下列条目与原文 §I.13 p.44 逐条对应；圆括号中的页码为原文页码。

\begin{description}[leftmargin=6em]
\item[$E$] 非阿基米德局部域，剩余域为 $\FF_q$（$q=p^f$，$p>0$）（p.44）
\item[$k$] $\FF_q$ 的代数闭包 $k=\ol{\FF}_q$（p.44）
\item[$\breve{E}$] $E$ 的完全无分歧扩的完备化，其剩余域为 $k$（p.44）
\item[$\ol{E}$] $E$ 的可分闭包（p.44）
\item[$\Gamma$] 绝对 Galois 群 $\Gamma=\Gal(\ol{E}|E)$（p.44）
\item[$W_E$] Weil 群（p.44）
\item[$I_E$] 惯性子群（p.44）
\item[$P_E$] 野惯性子群（p.44）
\item[$P$] 通常指 $P_E$ 的开子群；偶尔也指有限 $p$-群，或更一般地指阶与 $\ell$ 互素的有限可解群（p.44）
\item[$G$] $E$ 上的约化群（一律假定连通）（p.44）
\item[$\ul{X}$] 把 $S$ 送到连续映射 $|S|\to X$ 的层（$S$ 取在相关测试范畴中，通常是完满空间）（p.44）
\item[$|S|$] 空间 $S$ 的底层拓扑空间（配合 $\ul{X}$ 使用）（p.44）
\item[anima] 即 [Lur09] 中的“空间”、$\infty$-群胚或同伦型；此词不译（p.44）
\item[{$A[X_1,\dots,X_n]$}] 多项式代数；animated $A$-代数在筛余极限下由它们自由生成（p.44）
\item[$\mathcal{C}^{BG}$] $\mathcal{C}$ 中 $G$-等变对象的（$\infty$-）范畴，即函子 $BG\to\mathrm{Cat}_\infty$ 的极限（p.44）
\item[$\ast/G$] 分类叠（即常记作 $BG$ 者）；用此写法是为避与 Borel 子群 $B$ 混淆（p.44）
\item[$\ast$] 单点，即 $\mathrm{Spd}\,\FF_q$（p.44）
\item[$B(G)$] Kottwitz 集（p.44）
\item[$B$] 亦表示 Borel 子群，故分类叠写作 $\ast/B$ 而非 $BB$（p.44）
\end{description}
